\section{Reading the focused reduction paper}
\label{notes:paper-reading-guide}
The accompanying paper, \emph{Speeding up Feynman integral reduction with Kira: exact identities and learned choices}, reorganizes the completed study around the reduction problem. It first defines the integrals, gives a tadpole example of integration by parts, and then explains the exact identities and learned decisions. Its main comparison is the completed serial experiment with 93 requests. The present notes retain the original experiments, models, intermediate comparisons, and failures.

The final serial comparison gives median complete-collection times of 45.091 seconds for the preceding application and 41.137 seconds for the final application, with all 837 results verified. These values include processing all requests in one call. The earlier 75-request result, including its 39.1\% cumulative reduction across the vacuum and reconstruction changes, uses a different collection and timing procedure. The paper states both controls explicitly and does not combine the percentages.

The distinction between total savings and learned savings also remains explicit. The earlier seed and equation models improve their recorded controlled workloads. The reconstruction-era model-only comparison has a ratio of 0.983 with interval $[0.951,1.003]$, so it does not resolve a reliable additional gain. The latest arithmetic selector and nonplanar equation filter fail validation and remain disabled. Native row selection remains in use. The initial parallel comparison retains its unresolved 20-versus-16-master basis difference; the later successful serial test does not replace that evidence.

The revised paper gives the analytic derivations, accepted physical domains, exact comparison conditions, and slower subgroups beside the measured results. Higher-loop products and connected bubble trees remain distinct from general connected scattering. This editorial revision introduces no new training, native reduction, or timing measurement.
